\subsubsection{Chinese Remainder Theorem}
\begin{lstlisting}[style=mystyle, language=C++]
// TC: O(N * log(max(m))) donde N = num. congruencias
//     (cada modpow cuesta O(log m[i]))
// SC: O(1) extra
// usa: resolver sistema de congruencias pairwise coprimes
#include <bits/stdc++.h>
using namespace std;

// returns smallest x >= 0 such that x == r[i] (mod m[i]) for all i
//     assumes m[i] are pairwise coprime
long long crt(const vector<long long>& r, const vector<long long>& m){ 
	long long x = 0;
	long long M = 1; // modular inverse
	for(int i=0 ; i<r.size() ; i++){
		long long diff = (r[i] - x) % m[i];
		if(diff < 0) diff += m[i];
		long long inv = modpow(M % m[i], m[i] - 2, m[i]);  // (a^b) % m[i]
		long long t = diff * inv % m[i];
		x = x + M * t;
		M *= m[i];
		x %= M;
		if(x < 0) x += M;
	}
	return x;
}

int main(){
	// x == 2 (mod 3), x == 3 (mod 5), x == 2 (mod 7) -> x = 23
	cout << crt({2, 3, 2}, {3, 5, 7}) << "\n";
	return 0;
}
\end{lstlisting}
